\documentclass[12pt,a4paper]{report}

% ------------------ PACKAGES ------------------
\usepackage[a4paper,left=20mm,right=10mm,top=20mm,bottom=20mm]{geometry}
\usepackage{setspace}
\usepackage{titlesec}
\usepackage{indentfirst}
\usepackage{fancyhdr}
\usepackage{graphicx}
\usepackage{caption}
\usepackage{amsmath}
\usepackage{chngcntr}
\usepackage{enumitem}
\usepackage{float}
\usepackage{booktabs}
\usepackage{longtable}
\usepackage{listings}
\usepackage{newtxtext,newtxmath}

% ------------------ SPACING ------------------
\onehalfspacing
\setlength{\parindent}{1.25cm}
\setlength{\parskip}{0pt}

% ------------------ PAGE NUMBERING ------------------
\pagestyle{fancy}
\fancyhf{}
\fancyfoot[C]{\thepage}
\renewcommand{\headrulewidth}{0pt}

% ------------------ CHAPTER FORMAT ------------------
\titleformat{\chapter}
{\bfseries\centering\fontsize{13pt}{13pt}\selectfont}
{\thechapter}
{1em}
{}
\titlespacing*{\chapter}{0pt}{0pt}{1\baselineskip}

\titleformat{\section}
{\bfseries\fontsize{12pt}{12pt}\selectfont}
{\thesection}
{1em}
{}

% ------------------ NUMBERING ------------------
\numberwithin{equation}{chapter}
\numberwithin{figure}{chapter}
\numberwithin{table}{chapter}

\captionsetup[figure]{labelfont=bf, justification=centering}
\captionsetup[table]{font=small}

% ------------------ ENUMERATION STYLE ------------------
\setlist[itemize]{label=-, leftmargin=1.25cm}
\setlist[enumerate,1]{label=\alph*)}
\setlist[enumerate,2]{label=\arabic*)}

% ------------------ CODE STYLE ------------------
\lstset{
    basicstyle=\ttfamily\fontsize{9pt}{9pt}\selectfont,
    breaklines=true,
    numbers=left,
    numberstyle=\tiny,
    frame=single,
    tabsize=2
}

% ------------------ DOCUMENT ------------------
\begin{document}

% ------------------ TITLE PAGE ------------------
\begin{titlepage}
\centering
\begin{figure}[htbp]
\centering
\IfFileExists{img/utm_logo.png}{\includegraphics[width=0.2\textwidth]{img/utm_logo.png}}{\vspace{2cm}}
\end{figure}
MINISTRY OF EDUCATION AND RESEARCH\\
OF THE REPUBLIC OF MOLDOVA\\
Technical University of Moldova\\
Faculty of Computers, Informatics and Microelectronics\\
Department of Software Engineering and Automation\\[2cm]

Vasilean Vitalie\\
Group FAF-241\\[1.5cm]

\textbf{\Large Report}\\
Laboratory Work No. 8\\
Operating Systems\\
Topic: Processes Talking Through Pipes (xv6)\\[2cm]

\begin{flushright}
Checked by:\\
univ. assist. / lect. Patricia Reițman \\
Department of Software Engineering and Automation\\
FCIM, UTM
\end{flushright}

\vspace{3cm}

Chisinau -- 2026

\end{titlepage}

\pagenumbering{arabic}
\setcounter{page}{2}

% ------------------ TABLE OF CONTENTS ------------------
\tableofcontents
\newpage

% ------------------ CHAPTER 1 ------------------
\chapter{INTRODUCTION AND GOAL}

\section{Introduction}
Inside the xv6 teaching operating system, standard C libraries such as \texttt{glibc}, tracing utilities like \texttt{strace}, and manual pages do not exist. The user-space execution environment interfaces with the underlying RISC-V hardware solely through a minimalist set of 22 fundamental kernel system calls. Process management and inter-process communication rely directly on primitives exposed in \texttt{user/user.h}: \texttt{fork}, \texttt{pipe}, \texttt{read}, \texttt{write}, \texttt{close}, \texttt{wait}, \texttt{exit}, and \texttt{getpid}. Understanding how unidirectional byte streams function at this bare interface provides deep insight into kernel resource accounting, descriptor table isolation, and concurrent process scheduling without user-space runtime abstractions.

\section{Lab Objective}
The primary goals of this laboratory work are:
\begin{itemize}
    \item implementing \texttt{pingpong}, a synchronization protocol where parent and child processes bounce a byte bidirectionally using two distinct pipes;
    \item writing \texttt{primes}, Doug McIlroy's concurrent prime sieve pipeline using chained processes connected via pipes;
    \item managing file descriptor hygiene under strict kernel constraints, preventing file descriptor leaks (\texttt{NOFILE}) and avoiding indefinite pipeline hangs caused by unclosed write descriptors;
    \item observing kernel process management limits (\texttt{NPROC}) and recursion constraints enforced by compiler flags (\texttt{-Werror});
    \item implementing the optional stretch utility \texttt{xargs} to read newline-delimited command arguments from standard input and dispatch them via \texttt{fork}, \texttt{exec}, and \texttt{wait}.
\end{itemize}

% ------------------ CHAPTER 2 ------------------
\chapter{BIDIRECTIONAL COMMUNICATION: PINGPONG}

\section{Implementation of \texttt{user/pingpong.c}}
In xv6 and standard UNIX semantics, a pipe is strictly unidirectional. Bidirectional communication requires two separate pipes: one for parent-to-child transmission (\texttt{p2c}) and another for child-to-parent transmission (\texttt{c2p}). The child reads a single byte from \texttt{p2c}, prints its PID along with the confirmation message, and forwards the byte back via \texttt{c2p} before exiting.

\begin{lstlisting}[language=C, caption={Source code of \texttt{user/pingpong.c}}]
#include "kernel/types.h"
#include "user/user.h"

int
main(int argc, char *argv[])
{
    int p2c[2]; // parent -> child: [0] read, [1] write
    int c2p[2]; // child -> parent: [0] read, [1] write
    char byte = 'x';

    if (pipe(p2c) < 0 || pipe(c2p) < 0) {
        printf("pipe failed\n");
        exit(1);
    }

    int pid = fork();
    if (pid < 0) {
        printf("fork failed\n");
        exit(1);
    }

    if (pid == 0) {
        // Child process: close unused ends
        close(p2c[1]);
        close(c2p[0]);

        read(p2c[0], &byte, 1);
        printf("%d: received ping\n", getpid());
        close(p2c[0]);

        write(c2p[1], &byte, 1);
        close(c2p[1]);

        exit(0);
    } else {
        // Parent process: close unused ends
        close(p2c[0]);
        close(c2p[1]);

        write(p2c[1], &byte, 1);
        close(p2c[1]);

        read(c2p[0], &byte, 1);
        printf("%d: received pong\n", getpid());
        close(c2p[0]);

        wait(0);
        exit(0);
    }
}
\end{lstlisting}

\section{Execution and Verification}
After registering \texttt{\$U/\_pingpong\textbackslash} inside the \texttt{UPROGS} section of the \texttt{Makefile}, the kernel was built and launched inside QEMU:
\begin{lstlisting}[language=bash]
$ make qemu
...
xv6 kernel is booting

hart 2 starting
hart 1 starting
init: starting sh
$ pingpong
4: received ping
3: received pong
$ 
\end{lstlisting}

\section{Analysis and Observations}
\begin{enumerate}
    \item \textbf{Requirement for Two Independent Pipes:}
    A UNIX pipe provides a single unidirectional byte stream. If both processes shared a single pipe descriptor pair, a severe race condition would emerge: the parent writing the byte into the pipe could immediately execute a subsequent \texttt{read()} and consume its own transmitted byte before the child process is scheduled by the kernel. Allocating two distinct channels guarantees isolation and prevents self-reads.
    \item \textbf{Implicit Synchronization via Blocking Reads:}
    The source code contains no explicit polling loops or conditional synchronization barriers. Synchronization is achieved naturally through the blocking semantics of the \texttt{read()} system call. When the parent invokes \texttt{read(c2p[0], ...)}, the calling process is transitioned into a sleeping state by the xv6 scheduler because \texttt{c2p} contains no data. The parent remains suspended until the child executes \texttt{write(c2p[1], ...)}, which wakes the sleeping parent process.
    \item \textbf{Process Identifier (PID) Allocation Order:}
    The child process outputs a larger PID (\texttt{4}) than the parent (\texttt{3}). The xv6 kernel maintains a global monotonic integer counter for process identifiers. The parent process is allocated a PID upon its instantiation by the initial shell; when it later calls \texttt{fork()}, the kernel allocates the next sequential integer to the newly created child.
\end{enumerate}

% ------------------ CHAPTER 3 ------------------
\chapter{CONCURRENT PRIME SIEVE: PRIMES}

\section{Implementation of \texttt{user/primes.c}}
Doug McIlroy's pipeline model calculates prime numbers by generating an autonomous pipeline of processes. The root process generates integers from 2 through 35 and pushes them into an initial pipe. Each downstream stage reads its first available integer, which is mathematically guaranteed to be prime. The stage prints the value, creates a new pipe, forks a successor stage, and filters remaining integers: values non-divisible by its assigned prime are forwarded, while multiples are dropped.

\begin{lstlisting}[language=C, caption={Source code of \texttt{user/primes.c}}]
#include "kernel/types.h"
#include "user/user.h"

void stage(int left_fd) __attribute__((noreturn));

void
stage(int left_fd)
{
    int p;
    int n;

    int bytes = read(left_fd, &p, sizeof(p));
    if (bytes <= 0) {
        close(left_fd);
        exit(0);
    }

    printf("prime %d\n", p);

    int right_pipe[2];
    if (pipe(right_pipe) < 0) {
        close(left_fd);
        exit(1);
    }

    int pid = fork();
    if (pid < 0) {
        close(left_fd);
        close(right_pipe[0]);
        close(right_pipe[1]);
        exit(1);
    }

    if (pid == 0) {
        // Child stage: consume from right_pipe[0]
        close(left_fd);
        close(right_pipe[1]);
        stage(right_pipe[0]);
    } else {
        // Current stage: forward non-divisible numbers to right_pipe[1]
        close(right_pipe[0]);

        while (read(left_fd, &n, sizeof(n)) > 0) {
            if (n % p != 0) {
                write(right_pipe[1], &n, sizeof(n));
            }
        }

        close(left_fd);
        close(right_pipe[1]); // Close write end to transmit EOF
        wait(0);
        exit(0);
    }
}

int
main(int argc, char *argv[])
{
    int p[2];
    if (pipe(p) < 0) {
        exit(1);
    }

    int pid = fork();
    if (pid < 0) {
        exit(1);
    }

    if (pid == 0) {
        close(p[1]);
        stage(p[0]);
    } else {
        close(p[0]);
        for (int i = 2; i <= 35; i++) {
            write(p[1], &i, sizeof(i));
        }
        close(p[1]);
        wait(0);
        exit(0);
    }
}
\end{lstlisting}

\section{Execution and Verification}
Adding \texttt{\$U/\_primes\textbackslash} to the \texttt{Makefile} and executing under QEMU yielded the sequence of primes from 2 to 31:
\begin{lstlisting}[language=bash]
$ primes
prime 2
prime 3
prime 5
prime 7
prime 11
prime 13
prime 17
prime 19
prime 23
prime 29
prime 31
$ 
\end{lstlisting}

\section{Analysis and Observations}
\begin{enumerate}
    \item \textbf{Process Count at Execution Milestone:}
    When \texttt{"prime 31"} is output to the terminal, exactly 12 processes exist concurrently:
    \begin{equation}
        N_{\text{active}} = N_{\text{main}} + N_{\text{primes}} = 1 + 11 = 12
    \end{equation}
    where $N_{\text{main}}$ is the initial generator process, and $N_{\text{primes}}$ corresponds to the 11 discovered primes in the range $[2, 35]$: 2, 3, 5, 7, 11, 13, 17, 19, 23, 29, and 31.
    \item \textbf{Upper Limit Scaling (Limit 280 Experiment):}
    Increasing the generation boundary from 35 to 280 tests kernel capacity limits. In the interval $[2, 280]$, there are 59 prime numbers. In xv6, \texttt{kernel/param.h} fixes \texttt{NPROC = 64} and \texttt{NOFILE = 16}. With \texttt{init}, \texttt{sh}, and the initial generator process active, spawning 59 sieve stages requires 62 process slots, approaching the system ceiling. If all unused file descriptors are meticulously closed, the pipeline successfully finishes on stock xv6. However, if file descriptors are leaked, processes hit \texttt{NOFILE} (16 file descriptors) long before process limits are reached, resulting in failed \texttt{pipe()} invocations and early termination.
    \item \textbf{Parallelism Characteristics:}
    This architecture constitutes a parallel pipeline sieve because each stage functions as an independent concurrent filtering worker connected asynchronously via memory-buffered pipes. Later stages can filter and process earlier values while upstream processes simultaneously generate, test, and transmit subsequent numbers through the pipeline.
\end{enumerate}

% ------------------ CHAPTER 4 ------------------
\chapter{BONUS IMPLEMENTATION: XARGS}

\section{Implementation of \texttt{user/xargs.c}}
The \texttt{xargs} utility reads input from standard input line by line and executes the target utility for each row, appending the parsed line as command-line arguments. In xv6, no environment \texttt{PATH} exists, necessitating direct binary invocation via \texttt{exec}.

\begin{lstlisting}[language=C, caption={Source code of \texttt{user/xargs.c}}]
#include "kernel/types.h"
#include "kernel/param.h"
#include "user/user.h"

int
main(int argc, char *argv[])
{
    char *args[MAXARG];
    char buf[512];
    char ch;
    int i;

    if (argc < 2) {
        printf("Usage: xargs <command> [args...]\n");
        exit(1);
    }

    // Populate initial command arguments
    for (i = 1; i < argc; i++) {
        args[i - 1] = argv[i];
    }

    int base_argc = argc - 1;
    int pos = 0;

    // Read characters from standard input (file descriptor 0)
    while (read(0, &ch, 1) > 0) {
        if (ch == '\n') {
            buf[pos] = '\0';
            args[base_argc] = buf;
            args[base_argc + 1] = 0;

            int pid = fork();
            if (pid < 0) {
                printf("xargs: fork failed\n");
                exit(1);
            }

            if (pid == 0) {
                exec(args[0], args);
                printf("xargs: exec %s failed\n", args[0]);
                exit(1);
            }

            wait(0);
            pos = 0;
        } else {
            if (pos < (int)sizeof(buf) - 1) {
                buf[pos++] = ch;
            }
        }
    }

    // Flush any trailing line missing a newline character
    if (pos > 0) {
        buf[pos] = '\0';
        args[base_argc] = buf;
        args[base_argc + 1] = 0;

        int pid = fork();
        if (pid == 0) {
            exec(args[0], args);
            exit(1);
        }
        wait(0);
    }

    exit(0);
}
\end{lstlisting}

\section{Execution and Verification}
The command execution flow was verified against the requirement specified in the lab manual:
\begin{lstlisting}[language=bash]
$ echo hello too | xargs echo bye
bye hello too
$ 
\end{lstlisting}

% ------------------ CONCLUSION ------------------
\chapter*{CONCLUSION}
\addcontentsline{toc}{chapter}{CONCLUSION}
This laboratory work explored inter-process communication and concurrent scheduling inside the xv6 operating system using bare kernel system calls. Implementing low-level pipe primitives directly without runtime library abstractions provided practical insight into operating system resource isolation and scheduling mechanics.

Key conclusions derived from the experiments include:
\begin{itemize}
    \item \textbf{Descriptor Table Isolation:} Pipes are unidirectional byte buffers. Full-duplex messaging between parent and child processes requires two separate pipes to prevent self-consumption and race conditions.
    \item \textbf{Blocking Synchronization and EOF Mechanics:} A reading process blocks automatically when an incoming pipe is empty. Crucially, an end-of-file condition (\texttt{EOF}) is only raised when every write end referencing the pipe across the entire system is closed; omitting closures leads directly to deadlocks.
    \item \textbf{Kernel Table Limits:} Concurrent pipeline architectures are heavily bounded by process table capacity (\texttt{NPROC}) and per-process descriptor limits (\texttt{NOFILE}). Rigorous file descriptor cleanup is essential for sustained system stability.
\end{itemize}

% ------------------ REFERENCES ------------------
\chapter*{REFERENCES}
\addcontentsline{toc}{chapter}{REFERENCES}

\noindent [1] Russ Cox, Frans Kaashoek, and Robert Morris, \textit{xv6: a simple, Unix-like teaching operating system}, MIT CSAIL, 2024.

\noindent [2] M. Doug McIlroy, \textit{A Research UNIX Reader: System Architecture, 1972-1989}, AT\&T Bell Laboratories Computing Science Technical Report 139, 1987.

\noindent [3] Remzi H. Arpaci-Dusseau and Andrea C. Arpaci-Dusseau, \textit{Operating Systems: Three Easy Pieces}, Arpaci-Dusseau Books, 2018.

\noindent [4] Department of Software Engineering and Automation, \textit{Operating Systems Lab 8: Processes Talking Through Pipes (xv6)}, FCIM, Technical University of Moldova, 2026.

\end{document}
